
\subsubsection{Reference Trace}

A central question for implementation is whether a section, given only the 
beam strains $\bm{e} = (\bm{\gamma}, \bm{\kappa})$, can reconstruct the full 
three-dimensional strain distribution $\bm{\varepsilon}(x, \bm{r})$ without 
knowing its axial position $x$, or incorporating auxiliary warping fields \(\bm{\alpha}(x)\).

The three-dimensional strain field at section $x$ depends on the following 
\emph{local parameters}:
\begin{equation}
\hat{\bm{e}}(x) := \begin{pmatrix} 
\epsilon_0(x) \\ \bm{\beta} \\ \hat{\varkappa} \\ \bm{\epsilon}_1(x) 
\end{pmatrix} \in \mathbb{R}^6,
\label{eq:local-params}
\end{equation}
where:
\begin{itemize}
\item $\epsilon_0(x) := a_\varepsilon - x\,\CentroidLocation \cdot \bm{b}_{\Section}$ 
is the axial strain at the centroid,
\item $\bm{\epsilon}_1(x) := \bm{a}_\Omega + x\,\bm{b}_{\Section}$ is the axial 
strain gradient (curvature contribution),
\item $\hat{\varkappa} := a_\vartheta + c_\vartheta$ is the total twist rate 
(including flexure-induced twist),
\item $\bm{\beta} := \bm{b}_{\Section}$ is the flexural shear loading parameter.
\end{itemize}

These parameters suffice for pointwise strain evaluation:
\begin{align}
\varepsilon(\bm{r}) &= \epsilon_0 + \bm{r} \cdot \bm{\epsilon}_1, 
\label{eq:axial-strain-local} \\
\bm{\varepsilon}_\gamma(\bm{r}) &= \hat{\varkappa}\,(\FrameBasis \times \bm{r} + \nabla\WarpTwistIesan) 
+ \big(\bm{W}_{(b)} + (\nabla\bm{\WarpShearIesan})^\top\big)\,\bm{\beta}.
\label{eq:shear-strain-local}
\end{align}

The local parameters are related to the Ieşan parameters by
\begin{equation}
\hat{\bm{e}}(x) = \mathbf{L}(x)\,\bm{p} = (\mathbf{L}_0 + x\,\mathbf{L}_1)\,\bm{p},
\label{eq:q-from-p}
\end{equation}
where, in the partition $\hat{\bm{e}} = (\epsilon_0, \bm{\beta}, \hat{\varkappa}, \bm{\epsilon}_1)$ 
and $\bm{p} = (a_\varepsilon, \bm{a}_\Omega, a_\vartheta, \bm{b}_{\Section})$:
\begin{equation}
\mathbf{L}_0 := \begin{pmatrix}
1 & \mathbf{0}^\top & 0 & \mathbf{0}^\top \\
\mathbf{0} & \mathbf{0} & \mathbf{0} & \oneS \\
0 & \mathbf{0}^\top & 1 & \CenterShearLocation^\top \\
\mathbf{0} & \oneS & \mathbf{0} & \mathbf{0}
\end{pmatrix}, \qquad
\mathbf{L}_1 := \begin{pmatrix}
0 & \mathbf{0}^\top & 0 & -\CentroidLocation^\top \\
\mathbf{0} & \mathbf{0} & \mathbf{0} & \mathbf{0} \\
0 & \mathbf{0}^\top & 0 & \mathbf{0}^\top \\
\mathbf{0} & \mathbf{0} & \mathbf{0} & \oneS
\end{pmatrix}.
\label{eq:L-matrices}
\end{equation}
where the shear center $\CenterShearLocation$ is defined by \eqref{eq:def-shear-center}, 
and the centroid location $\CentroidLocation$ by \eqref{eq:def-centroid}.

\subsubsection{Implementable Traces}

A trace is \emph{implementable} if a section can recover the local parameters 
$\hat{\bm{e}}$ from the beam strains $\bm{e}$ without knowledge of $x$. This requires 
the existence of an $x$-independent matrix $\mathbf{X}$ such that
\begin{equation}
\bm{e}(x) = \mathbf{X}\,\hat{\bm{e}}(x)
\label{eq:implementable-def}
\end{equation}
for all Saint-Venant solutions. 
%
Substituting \eqref{eq:strain-polynomial} and \eqref{eq:q-from-p}:
\begin{equation}
\sum_m x^m\,\mathbf{T}_m\,\bm{p} = \mathbf{X}\,(\mathbf{L}_0 + x\,\mathbf{L}_1)\,\bm{p}.
\end{equation}
Matching powers of $x$ yields:
\begin{subequations}
\begin{align}
\TraceDeDpRoot &= \mathbf{X}\,\mathbf{L}_0, \label{eq:T0-condition} \\
\TraceDeDpOne &= \mathbf{X}\,\mathbf{L}_1, \label{eq:T1-condition} \\
\mathbf{T}_m &= \mathbf{0} \quad \text{for } m \geq 2. \label{eq:Tm-condition}
\end{align}
\end{subequations}

Since $\mathbf{L}_0$ is invertible, $\mathbf{X} = \TraceDeDpRoot\,\mathbf{L}_0^{-1}$, 
and the consistency condition becomes:
\begin{equation}
\boxed{\TraceDeDpOne = \TraceDeDpRoot\,\mathbf{L}_0^{-1}\,\mathbf{L}_1}
\label{eq:implementability-condition}
\end{equation}
where
$$
\mathbf{L}_0^{-1} = \begin{pmatrix}
1 & \mathbf{0}^\top & 0 & \mathbf{0}^\top \\
\mathbf{0} & \mathbf{0} & \mathbf{0} & \oneS \\
0 & -\CenterShearLocation^\top & 1 & \mathbf{0}^\top \\
\mathbf{0} & \oneS & \mathbf{0} & \mathbf{0}
\end{pmatrix}
\qquad \text{and}\qquad
\mathbf{L}_0^{-1}\mathbf{L}_1 = \begin{pmatrix}
0 & \mathbf{0}^\top & 0 & -\CentroidLocation^\top \\
\mathbf{0} & \mathbf{0} & \mathbf{0} & \oneS\\
0 & \mathbf{0}^\top & 0 & \mathbf{0}^\top \\
\mathbf{0} & \mathbf{0} & \mathbf{0} & \mathbf{0} \\
\end{pmatrix}$$

This condition constrains how the $x$-dependence in the beam strains must 
arise: it must come precisely from the $x$-dependence of the local parameters 
$(\epsilon_0(x), \bm{\epsilon}_1(x))$, not from any independent geometric coupling.

The space of implementable traces is parameterized by $GL(6)$: given any 
invertible $\mathbf{X}$, defining $\TraceDeDpRoot := \mathbf{X}\,\mathbf{L}_0$ 
and $\TraceDeDpOne := \mathbf{X}\,\mathbf{L}_1$ automatically satisfies 
\eqref{eq:implementability-condition}. The matrix $\mathbf{X}$ represents a 
linear relabeling of the local parameters as beam strains.

For an implementable trace, the section implementation is straightforward. 
Given beam strains $\bm{e}$, the section computes
\begin{equation}
\hat{\bm{e}} = \mathbf{X}^{-1}\,\bm{e} = \mathbf{L}_0\,\TraceDpDeRoot\,\bm{e}
\label{eq:q-recovery}
\end{equation}
and evaluates the three-dimensional strains at each fiber via 
\eqref{eq:axial-strain-local}--\eqref{eq:shear-strain-local}.
